% --- Archon physics formalization source begin ---
% archon:physics
% archon:covers PhyXMiniProblems/problem_phyx_mini_0730.lean
% archon:source-report reports/phyx_mini/problem_phyx_mini_0730.source.json

\chapter{Physics problem phyx\_mini\_0730}
\label{ch:PhyXMiniProblems_problem_phyx_mini_0730}

\paragraph{Problem source.}
\#\# Physical scenario

When a high-speed passenger train traveling at \$161\textbackslash{} km/h\$ rounds a bend, the engineer is shocked to see that a locomotive has improperly entered onto the track from a siding and is a distance \$D = 676\textbackslash{} m\$ ahead. The locomotive is moving at \$29.0\textbackslash{} km/h\$. The engineer of the high-speed train immediately applies the brakes.

\#\# Question

What must be the magnitude of the resulting constant deceleration if a collision is to be just avoided?

\#\# Answer choices

- A: A: \$0.824\textbackslash{} m/s\textasciicircum{}2\$

- B: B: \$0.886\textbackslash{} m/s\textasciicircum{}2\$

- C: C: \$0.994\textbackslash{} m/s\textasciicircum{}2\$

- D: D: \$0.998\textbackslash{} m/s\textasciicircum{}2\$

\#\# Figure caption (auxiliary; use the image as primary evidence)

The image depicts a diagram featuring two main elements on a railway: a high-speed train and a locomotive. 

1. **High-Speed Train**: 
   - Located on the left, comprised of several interconnected carriages.
   - Positioned on a curved section of track.
   - Labeled "High-speed train."
   - An arrow pointing to the right indicates its direction.

2. **Locomotive**:
   - Located on the right, depicted as a singular unit.
   - Positioned on a straight section of track before a switch point where the tracks diverge.
   - Labeled "Locomotive."
   - An arrow pointing to the right indicates its direction.

3. **Railway Tracks**:
   - Two sets of tracks: one for the high-speed train and one that diverges leading to the locomotive.
   - The tracks are illustrated with typical railway ties.

4. **Label "D"**:
   - Positioned between the high-speed train and the locomotive, indicating a distance.
   - Arrowed lines at each end of the "D" demonstrate the measurement of this distance between the two elements.

The overall layout shows a potential collision scenario or a merging path setup for illustrative or educational purposes.

\#\# Dataset metadata

Dynamics; Physical Model Grounding Reasoning, Multi-Formula Reasoning

\paragraph{Recorded answer/context.}
C: C: \$0.994\textbackslash{} m/s\textasciicircum{}2\$

\paragraph{Figure/image path.}
/root/proposal\_for\_physic/science-mango/phyx\_mini\_run/phyx\_data/test\_image/730.png

\paragraph{Formalization target.}
create a compiling Lean file with sorry bodies at `PhyXMiniProblems/problem\_phyx\_mini\_0730.lean`. The Lean declarations must preserve the physical quantities, dimensions or dimensional roles, figure labels, governing-law hypotheses, and final relation expressed by this problem.
Use Mathlib/Physlib names found through LeanExplore where available. If a physics API is missing, introduce faithful local abstractions rather than scalar placeholder aliases.

\begin{theorem}[Physics formalization target]
\label{thm:physics:phyx_mini_0730:target}
\lean{PhyXMiniProblems.ProblemPhyXMini0730.problem_phyx_mini_0730}
\uses{def:physics:phyx-mini-0730:phyxminiproblems-problemphyxmini0730-accelerationinmeterspersecondsquared, def:physics:phyx-mini-0730:phyxminiproblems-problemphyxmini0730-trainbrakingsetup, def:physics:phyx-mini-0730:phyxminiproblems-problemphyxmini0730-matchessuppliedrailwayfigure, def:physics:phyx-mini-0730:phyxminiproblems-problemphyxmini0730-matchesproblemreadouts, def:physics:phyx-mini-0730:phyxminiproblems-problemphyxmini0730-hasphysicaltrainbrakingparameters, def:physics:phyx-mini-0730:phyxminiproblems-problemphyxmini0730-satisfiestrainkinematics, def:physics:phyx-mini-0730:phyxminiproblems-problemphyxmini0730-satisfieslimitingnocollisioncondition, def:physics:phyx-mini-0730:phyxminiproblems-problemphyxmini0730-matchesdisplayeddecelerationchoice}
The assigned autoformalize agent should translate this physics problem into Lean declarations in the covered file, with theorem and lemma proof bodies written as `by sorry`.
\end{theorem}
\begin{proof}
This is an autoformalization task, not a proof task. Produce statements that can later be proved by the physics prover without weakening the physical model.
\end{proof}

% --- Archon declaration topology begin ---
\section{Declaration topology}
\begin{definition}[acceleration Dimension]
  \label{def:physics:phyx-mini-0730:phyxminiproblems-problemphyxmini0730-accelerationdimension}
  \lean{PhyXMiniProblems.ProblemPhyXMini0730.accelerationDimension}
  The physical dimension of longitudinal acceleration, L T⁻²
\end{definition}
\begin{proof}
  This is the declared model definition of acceleration Dimension.
\end{proof}
\begin{definition}[Length Magnitude]
  \label{def:physics:phyx-mini-0730:phyxminiproblems-problemphyxmini0730-lengthmagnitude}
  \lean{PhyXMiniProblems.ProblemPhyXMini0730.LengthMagnitude}
  A nonnegative physical length magnitude
\end{definition}
\begin{proof}
  This is the declared model definition of Length Magnitude.
\end{proof}
\begin{definition}[Longitudinal Position]
  \label{def:physics:phyx-mini-0730:phyxminiproblems-problemphyxmini0730-longitudinalposition}
  \lean{PhyXMiniProblems.ProblemPhyXMini0730.LongitudinalPosition}
  A signed longitudinal position along the common track direction
\end{definition}
\begin{proof}
  This is the declared model definition of Longitudinal Position.
\end{proof}
\begin{definition}[Time Duration]
  \label{def:physics:phyx-mini-0730:phyxminiproblems-problemphyxmini0730-timeduration}
  \lean{PhyXMiniProblems.ProblemPhyXMini0730.TimeDuration}
  A nonnegative elapsed-time duration
\end{definition}
\begin{proof}
  This is the declared model definition of Time Duration.
\end{proof}
\begin{definition}[Speed Magnitude]
  \label{def:physics:phyx-mini-0730:phyxminiproblems-problemphyxmini0730-speedmagnitude}
  \lean{PhyXMiniProblems.ProblemPhyXMini0730.SpeedMagnitude}
  A nonnegative physical speed magnitude
\end{definition}
\begin{proof}
  This is the declared model definition of Speed Magnitude.
\end{proof}
\begin{definition}[Signed Speed Quantity]
  \label{def:physics:phyx-mini-0730:phyxminiproblems-problemphyxmini0730-signedspeedquantity}
  \lean{PhyXMiniProblems.ProblemPhyXMini0730.SignedSpeedQuantity}
  A signed longitudinal velocity component
\end{definition}
\begin{proof}
  This is the declared model definition of Signed Speed Quantity.
\end{proof}
\begin{definition}[Acceleration Magnitude]
  \label{def:physics:phyx-mini-0730:phyxminiproblems-problemphyxmini0730-accelerationmagnitude}
  \lean{PhyXMiniProblems.ProblemPhyXMini0730.AccelerationMagnitude}
  \uses{def:physics:phyx-mini-0730:phyxminiproblems-problemphyxmini0730-accelerationdimension}
  A nonnegative magnitude of longitudinal acceleration
\end{definition}
\begin{proof}
  This is the declared model definition of Acceleration Magnitude.
\end{proof}
\begin{definition}[length Readout]
  \label{def:physics:phyx-mini-0730:phyxminiproblems-problemphyxmini0730-lengthreadout}
  \lean{PhyXMiniProblems.ProblemPhyXMini0730.lengthReadout}
  \uses{def:physics:phyx-mini-0730:phyxminiproblems-problemphyxmini0730-lengthmagnitude}
  Read a physical length magnitude in a selected length unit
\end{definition}
\begin{proof}
  This is the declared model definition of length Readout.
\end{proof}
\begin{definition}[position Readout]
  \label{def:physics:phyx-mini-0730:phyxminiproblems-problemphyxmini0730-positionreadout}
  \lean{PhyXMiniProblems.ProblemPhyXMini0730.positionReadout}
  \uses{def:physics:phyx-mini-0730:phyxminiproblems-problemphyxmini0730-longitudinalposition}
  Read a signed longitudinal position in a selected length unit
\end{definition}
\begin{proof}
  This is the declared model definition of position Readout.
\end{proof}
\begin{definition}[time Readout]
  \label{def:physics:phyx-mini-0730:phyxminiproblems-problemphyxmini0730-timereadout}
  \lean{PhyXMiniProblems.ProblemPhyXMini0730.timeReadout}
  \uses{def:physics:phyx-mini-0730:phyxminiproblems-problemphyxmini0730-timeduration}
  Read an elapsed time in a selected time unit
\end{definition}
\begin{proof}
  This is the declared model definition of time Readout.
\end{proof}
\begin{definition}[speed Readout]
  \label{def:physics:phyx-mini-0730:phyxminiproblems-problemphyxmini0730-speedreadout}
  \lean{PhyXMiniProblems.ProblemPhyXMini0730.speedReadout}
  \uses{def:physics:phyx-mini-0730:phyxminiproblems-problemphyxmini0730-speedmagnitude}
  Read a speed magnitude in coherent selected length and time units
\end{definition}
\begin{proof}
  This is the declared model definition of speed Readout.
\end{proof}
\begin{definition}[signed Speed Readout]
  \label{def:physics:phyx-mini-0730:phyxminiproblems-problemphyxmini0730-signedspeedreadout}
  \lean{PhyXMiniProblems.ProblemPhyXMini0730.signedSpeedReadout}
  \uses{def:physics:phyx-mini-0730:phyxminiproblems-problemphyxmini0730-signedspeedquantity}
  Read a signed longitudinal velocity in selected units
\end{definition}
\begin{proof}
  This is the declared model definition of signed Speed Readout.
\end{proof}
\begin{definition}[acceleration Readout]
  \label{def:physics:phyx-mini-0730:phyxminiproblems-problemphyxmini0730-accelerationreadout}
  \lean{PhyXMiniProblems.ProblemPhyXMini0730.accelerationReadout}
  \uses{def:physics:phyx-mini-0730:phyxminiproblems-problemphyxmini0730-accelerationmagnitude}
  Read an acceleration magnitude in coherent selected units
\end{definition}
\begin{proof}
  This is the declared model definition of acceleration Readout.
\end{proof}
\begin{definition}[length In Meters]
  \label{def:physics:phyx-mini-0730:phyxminiproblems-problemphyxmini0730-lengthinmeters}
  \lean{PhyXMiniProblems.ProblemPhyXMini0730.lengthInMeters}
  \uses{def:physics:phyx-mini-0730:phyxminiproblems-problemphyxmini0730-lengthmagnitude, def:physics:phyx-mini-0730:phyxminiproblems-problemphyxmini0730-lengthreadout}
  Metre readout for physical lengths
\end{definition}
\begin{proof}
  This is the declared model definition of length In Meters.
\end{proof}
\begin{definition}[position In Meters]
  \label{def:physics:phyx-mini-0730:phyxminiproblems-problemphyxmini0730-positioninmeters}
  \lean{PhyXMiniProblems.ProblemPhyXMini0730.positionInMeters}
  \uses{def:physics:phyx-mini-0730:phyxminiproblems-problemphyxmini0730-longitudinalposition, def:physics:phyx-mini-0730:phyxminiproblems-problemphyxmini0730-positionreadout}
  Metre coordinate along the right-positive track direction
\end{definition}
\begin{proof}
  This is the declared model definition of position In Meters.
\end{proof}
\begin{definition}[time In Seconds]
  \label{def:physics:phyx-mini-0730:phyxminiproblems-problemphyxmini0730-timeinseconds}
  \lean{PhyXMiniProblems.ProblemPhyXMini0730.timeInSeconds}
  \uses{def:physics:phyx-mini-0730:phyxminiproblems-problemphyxmini0730-timeduration, def:physics:phyx-mini-0730:phyxminiproblems-problemphyxmini0730-timereadout}
  Second readout for the limiting time
\end{definition}
\begin{proof}
  This is the declared model definition of time In Seconds.
\end{proof}
\begin{definition}[speed In Kilometers Per Hour]
  \label{def:physics:phyx-mini-0730:phyxminiproblems-problemphyxmini0730-speedinkilometersperhour}
  \lean{PhyXMiniProblems.ProblemPhyXMini0730.speedInKilometersPerHour}
  \uses{def:physics:phyx-mini-0730:phyxminiproblems-problemphyxmini0730-speedmagnitude, def:physics:phyx-mini-0730:phyxminiproblems-problemphyxmini0730-speedreadout}
  Kilometre-per-hour readout for the two stated initial speeds
\end{definition}
\begin{proof}
  This is the declared model definition of speed In Kilometers Per Hour.
\end{proof}
\begin{definition}[speed In Meters Per Second]
  \label{def:physics:phyx-mini-0730:phyxminiproblems-problemphyxmini0730-speedinmeterspersecond}
  \lean{PhyXMiniProblems.ProblemPhyXMini0730.speedInMetersPerSecond}
  \uses{def:physics:phyx-mini-0730:phyxminiproblems-problemphyxmini0730-speedmagnitude, def:physics:phyx-mini-0730:phyxminiproblems-problemphyxmini0730-speedreadout}
  Metre-per-second readout for a speed magnitude
\end{definition}
\begin{proof}
  This is the declared model definition of speed In Meters Per Second.
\end{proof}
\begin{definition}[velocity In Meters Per Second]
  \label{def:physics:phyx-mini-0730:phyxminiproblems-problemphyxmini0730-velocityinmeterspersecond}
  \lean{PhyXMiniProblems.ProblemPhyXMini0730.velocityInMetersPerSecond}
  \uses{def:physics:phyx-mini-0730:phyxminiproblems-problemphyxmini0730-signedspeedquantity, def:physics:phyx-mini-0730:phyxminiproblems-problemphyxmini0730-signedspeedreadout}
  Signed metre-per-second longitudinal velocity component
\end{definition}
\begin{proof}
  This is the declared model definition of velocity In Meters Per Second.
\end{proof}
\begin{definition}[acceleration In Meters Per Second Squared]
  \label{def:physics:phyx-mini-0730:phyxminiproblems-problemphyxmini0730-accelerationinmeterspersecondsquared}
  \lean{PhyXMiniProblems.ProblemPhyXMini0730.accelerationInMetersPerSecondSquared}
  \uses{def:physics:phyx-mini-0730:phyxminiproblems-problemphyxmini0730-accelerationmagnitude, def:physics:phyx-mini-0730:phyxminiproblems-problemphyxmini0730-accelerationreadout}
  Metre-per-second-squared readout for the braking magnitude
\end{definition}
\begin{proof}
  This is the declared model definition of acceleration In Meters Per Second Squared.
\end{proof}
\begin{definition}[Rail Vehicle]
  \label{def:physics:phyx-mini-0730:phyxminiproblems-problemphyxmini0730-railvehicle}
  \lean{PhyXMiniProblems.ProblemPhyXMini0730.RailVehicle}
  The two rail vehicles labelled in image 730.png
\end{definition}
\begin{proof}
  This is the declared model definition of Rail Vehicle.
\end{proof}
\begin{definition}[Travel Direction]
  \label{def:physics:phyx-mini-0730:phyxminiproblems-problemphyxmini0730-traveldirection}
  \lean{PhyXMiniProblems.ProblemPhyXMini0730.TravelDirection}
  Longitudinal direction of each blue velocity arrow in the image
\end{definition}
\begin{proof}
  This is the declared model definition of Travel Direction.
\end{proof}
\begin{definition}[Track Region]
  \label{def:physics:phyx-mini-0730:phyxminiproblems-problemphyxmini0730-trackregion}
  \lean{PhyXMiniProblems.ProblemPhyXMini0730.TrackRegion}
  The three track portions visible in the primary image
\end{definition}
\begin{proof}
  This is the declared model definition of Track Region.
\end{proof}
\begin{definition}[Figure Text Label]
  \label{def:physics:phyx-mini-0730:phyxminiproblems-problemphyxmini0730-figuretextlabel}
  \lean{PhyXMiniProblems.ProblemPhyXMini0730.FigureTextLabel}
  Literal text labels visible in the primary image
\end{definition}
\begin{proof}
  This is the declared model definition of Figure Text Label.
\end{proof}
\begin{definition}[Distance Marker Endpoint]
  \label{def:physics:phyx-mini-0730:phyxminiproblems-problemphyxmini0730-distancemarkerendpoint}
  \lean{PhyXMiniProblems.ProblemPhyXMini0730.DistanceMarkerEndpoint}
  Endpoints of the double-headed distance marker labelled D
\end{definition}
\begin{proof}
  This is the declared model definition of Distance Marker Endpoint.
\end{proof}
\begin{definition}[Supplied Railway Figure]
  \label{def:physics:phyx-mini-0730:phyxminiproblems-problemphyxmini0730-suppliedrailwayfigure}
  \lean{PhyXMiniProblems.ProblemPhyXMini0730.SuppliedRailwayFigure}
  \uses{def:physics:phyx-mini-0730:phyxminiproblems-problemphyxmini0730-railvehicle, def:physics:phyx-mini-0730:phyxminiproblems-problemphyxmini0730-traveldirection, def:physics:phyx-mini-0730:phyxminiproblems-problemphyxmini0730-trackregion, def:physics:phyx-mini-0730:phyxminiproblems-problemphyxmini0730-figuretextlabel, def:physics:phyx-mini-0730:phyxminiproblems-problemphyxmini0730-distancemarkerendpoint}
  Qualitative and numerical information transcribed from the supplied figure. The scalar marker readout is figure evidence; a separate premise relates it to the independent physical initial separation
\end{definition}
\begin{proof}
  This is the declared model definition of Supplied Railway Figure.
\end{proof}
\begin{definition}[Train Braking Setup]
  \label{def:physics:phyx-mini-0730:phyxminiproblems-problemphyxmini0730-trainbrakingsetup}
  \lean{PhyXMiniProblems.ProblemPhyXMini0730.TrainBrakingSetup}
  \uses{def:physics:phyx-mini-0730:phyxminiproblems-problemphyxmini0730-lengthmagnitude, def:physics:phyx-mini-0730:phyxminiproblems-problemphyxmini0730-longitudinalposition, def:physics:phyx-mini-0730:phyxminiproblems-problemphyxmini0730-timeduration, def:physics:phyx-mini-0730:phyxminiproblems-problemphyxmini0730-speedmagnitude, def:physics:phyx-mini-0730:phyxminiproblems-problemphyxmini0730-signedspeedquantity, def:physics:phyx-mini-0730:phyxminiproblems-problemphyxmini0730-accelerationmagnitude, def:physics:phyx-mini-0730:phyxminiproblems-problemphyxmini0730-railvehicle, def:physics:phyx-mini-0730:phyxminiproblems-problemphyxmini0730-traveldirection, def:physics:phyx-mini-0730:phyxminiproblems-problemphyxmini0730-suppliedrailwayfigure}
  Independent physical quantities and trajectories in the scenario. Coordinate time is measured in seconds from the instant braking begins. In particular, the braking deceleration is an independent physical quantity, not a definition of the recorded answer
\end{definition}
\begin{proof}
  This is the declared model definition of Train Braking Setup.
\end{proof}
\begin{definition}[Matches Supplied Railway Figure]
  \label{def:physics:phyx-mini-0730:phyxminiproblems-problemphyxmini0730-matchessuppliedrailwayfigure}
  \lean{PhyXMiniProblems.ProblemPhyXMini0730.MatchesSuppliedRailwayFigure}
  \uses{def:physics:phyx-mini-0730:phyxminiproblems-problemphyxmini0730-lengthinmeters, def:physics:phyx-mini-0730:phyxminiproblems-problemphyxmini0730-trainbrakingsetup}
  Primary-image evidence: the high-speed train is on the curved main-line approach, the locomotive is on the straight main line beyond the siding merge, both arrows point right, and D spans the high-speed train's front and the locomotive's rear
\end{definition}
\begin{proof}
  This is the declared model definition of Matches Supplied Railway Figure.
\end{proof}
\begin{definition}[Matches Problem Readouts]
  \label{def:physics:phyx-mini-0730:phyxminiproblems-problemphyxmini0730-matchesproblemreadouts}
  \lean{PhyXMiniProblems.ProblemPhyXMini0730.MatchesProblemReadouts}
  \uses{def:physics:phyx-mini-0730:phyxminiproblems-problemphyxmini0730-lengthinmeters, def:physics:phyx-mini-0730:phyxminiproblems-problemphyxmini0730-speedinkilometersperhour, def:physics:phyx-mini-0730:phyxminiproblems-problemphyxmini0730-trainbrakingsetup}
  Numerical data and direction conventions stated in the problem. The deceleration magnitude and limiting time have no numerical readout here
\end{definition}
\begin{proof}
  This is the declared model definition of Matches Problem Readouts.
\end{proof}
\begin{definition}[Has Physical Train Braking Parameters]
  \label{def:physics:phyx-mini-0730:phyxminiproblems-problemphyxmini0730-hasphysicaltrainbrakingparameters}
  \lean{PhyXMiniProblems.ProblemPhyXMini0730.HasPhysicalTrainBrakingParameters}
  \uses{def:physics:phyx-mini-0730:phyxminiproblems-problemphyxmini0730-lengthinmeters, def:physics:phyx-mini-0730:phyxminiproblems-problemphyxmini0730-speedinmeterspersecond, def:physics:phyx-mini-0730:phyxminiproblems-problemphyxmini0730-accelerationinmeterspersecondsquared, def:physics:phyx-mini-0730:phyxminiproblems-problemphyxmini0730-trainbrakingsetup}
  Positivity and ordering conditions for the physical braking regime
\end{definition}
\begin{proof}
  This is the declared model definition of Has Physical Train Braking Parameters.
\end{proof}
\begin{definition}[Satisfies Train Kinematics]
  \label{def:physics:phyx-mini-0730:phyxminiproblems-problemphyxmini0730-satisfiestrainkinematics}
  \lean{PhyXMiniProblems.ProblemPhyXMini0730.SatisfiesTrainKinematics}
  \uses{def:physics:phyx-mini-0730:phyxminiproblems-problemphyxmini0730-lengthinmeters, def:physics:phyx-mini-0730:phyxminiproblems-problemphyxmini0730-positioninmeters, def:physics:phyx-mini-0730:phyxminiproblems-problemphyxmini0730-speedinmeterspersecond, def:physics:phyx-mini-0730:phyxminiproblems-problemphyxmini0730-velocityinmeterspersecond, def:physics:phyx-mini-0730:phyxminiproblems-problemphyxmini0730-accelerationinmeterspersecondsquared, def:physics:phyx-mini-0730:phyxminiproblems-problemphyxmini0730-trainbrakingsetup}
  The high-speed train has constant rightward acceleration component -a, while the locomotive continues at its constant rightward speed. These are general one-dimensional kinematic laws. They contain neither an answer-choice value nor a closed formula for the requested acceleration
\end{definition}
\begin{proof}
  This is the declared model definition of Satisfies Train Kinematics.
\end{proof}
\begin{definition}[Satisfies Limiting No Collision Condition]
  \label{def:physics:phyx-mini-0730:phyxminiproblems-problemphyxmini0730-satisfieslimitingnocollisioncondition}
  \lean{PhyXMiniProblems.ProblemPhyXMini0730.SatisfiesLimitingNoCollisionCondition}
  \uses{def:physics:phyx-mini-0730:phyxminiproblems-problemphyxmini0730-positioninmeters, def:physics:phyx-mini-0730:phyxminiproblems-problemphyxmini0730-timeinseconds, def:physics:phyx-mini-0730:phyxminiproblems-problemphyxmini0730-velocityinmeterspersecond, def:physics:phyx-mini-0730:phyxminiproblems-problemphyxmini0730-trainbrakingsetup}
  The independent limiting time interprets “a collision is just avoided”: the front of the following train remains strictly behind the locomotive's rear at all earlier nonnegative times, and at the limiting zero-gap instant their positions and velocities agree. This is a boundary-event condition, not the requested numerical deceleration
\end{definition}
\begin{proof}
  This is the declared model definition of Satisfies Limiting No Collision Condition.
\end{proof}
\begin{definition}[Answer Choice]
  \label{def:physics:phyx-mini-0730:phyxminiproblems-problemphyxmini0730-answerchoice}
  \lean{PhyXMiniProblems.ProblemPhyXMini0730.AnswerChoice}
  Labels of the four displayed deceleration choices
\end{definition}
\begin{proof}
  This is the declared model definition of Answer Choice.
\end{proof}
\begin{definition}[displayed Deceleration In Meters Per Second Squared]
  \label{def:physics:phyx-mini-0730:phyxminiproblems-problemphyxmini0730-displayeddecelerationinmeterspersecondsquared}
  \lean{PhyXMiniProblems.ProblemPhyXMini0730.displayedDecelerationInMetersPerSecondSquared}
  \uses{def:physics:phyx-mini-0730:phyxminiproblems-problemphyxmini0730-answerchoice}
  Acceleration value in m/s² printed beside each answer label
\end{definition}
\begin{proof}
  This is the declared model definition of displayed Deceleration In Meters Per Second Squared.
\end{proof}
\begin{definition}[recorded Dataset Answer]
  \label{def:physics:phyx-mini-0730:phyxminiproblems-problemphyxmini0730-recordeddatasetanswer}
  \lean{PhyXMiniProblems.ProblemPhyXMini0730.recordedDatasetAnswer}
  \uses{def:physics:phyx-mini-0730:phyxminiproblems-problemphyxmini0730-answerchoice}
  The answer label recorded by the source dataset, retained as metadata
\end{definition}
\begin{proof}
  This is the declared model definition of recorded Dataset Answer.
\end{proof}
\begin{definition}[Rounds To Nearest Thousandth]
  \label{def:physics:phyx-mini-0730:phyxminiproblems-problemphyxmini0730-roundstonearestthousandth}
  \lean{PhyXMiniProblems.ProblemPhyXMini0730.RoundsToNearestThousandth}
  Agreement after rounding an acceleration to three decimal places
\end{definition}
\begin{proof}
  This is the declared model definition of Rounds To Nearest Thousandth.
\end{proof}
\begin{definition}[Matches Displayed Deceleration Choice]
  \label{def:physics:phyx-mini-0730:phyxminiproblems-problemphyxmini0730-matchesdisplayeddecelerationchoice}
  \lean{PhyXMiniProblems.ProblemPhyXMini0730.MatchesDisplayedDecelerationChoice}
  \uses{def:physics:phyx-mini-0730:phyxminiproblems-problemphyxmini0730-accelerationinmeterspersecondsquared, def:physics:phyx-mini-0730:phyxminiproblems-problemphyxmini0730-trainbrakingsetup, def:physics:phyx-mini-0730:phyxminiproblems-problemphyxmini0730-answerchoice, def:physics:phyx-mini-0730:phyxminiproblems-problemphyxmini0730-displayeddecelerationinmeterspersecondsquared, def:physics:phyx-mini-0730:phyxminiproblems-problemphyxmini0730-roundstonearestthousandth}
  The physical braking magnitude agrees with a displayed choice
\end{definition}
\begin{proof}
  This is the declared model definition of Matches Displayed Deceleration Choice.
\end{proof}
\begin{lemma}[braking Deceleration In Meters Per Second Squared exact]
  \label{lem:physics:phyx-mini-0730:phyxminiproblems-problemphyxmini0730-brakingdecelerationinmeterspersecondsquared-exact}
  \lean{PhyXMiniProblems.ProblemPhyXMini0730.brakingDecelerationInMetersPerSecondSquared_exact}
  \uses{def:physics:phyx-mini-0730:phyxminiproblems-problemphyxmini0730-accelerationinmeterspersecondsquared, def:physics:phyx-mini-0730:phyxminiproblems-problemphyxmini0730-trainbrakingsetup, def:physics:phyx-mini-0730:phyxminiproblems-problemphyxmini0730-matchessuppliedrailwayfigure, def:physics:phyx-mini-0730:phyxminiproblems-problemphyxmini0730-matchesproblemreadouts, def:physics:phyx-mini-0730:phyxminiproblems-problemphyxmini0730-hasphysicaltrainbrakingparameters, def:physics:phyx-mini-0730:phyxminiproblems-problemphyxmini0730-satisfiestrainkinematics, def:physics:phyx-mini-0730:phyxminiproblems-problemphyxmini0730-satisfieslimitingnocollisioncondition}
  The relative-motion kinematics determine the exact limiting deceleration as (161 - 29)\textasciicircum{}2 / (2 * 676 * 3.6\textasciicircum{}2) = 3025/3042 m/s²
\end{lemma}
\begin{proof}
  The conclusion follows from the governing relations and figure readouts named in the statement of braking Deceleration In Meters Per Second Squared exact.
\end{proof}
% --- Archon declaration topology end ---
% --- Archon physics formalization source end ---
